\chapter{Hamiltonian Mechanics And Phase Space}
\label{chap:hamiltonian-phase-space}

The Lagrangian chapter described motion as stationary paths in configuration
space. Hamiltonian mechanics reorganizes the same dynamics on phase space, the
space whose points specify both position and momentum. This reorganization is
not merely a change of variables. It makes the symplectic form, Poisson
brackets, canonical transformations, and phase-space volume available as
structural tools. These tools are the language in which integrability,
action-angle variables, perturbation theory, and KAM theory are most naturally
stated.

Throughout this chapter \(Q\) is a smooth configuration manifold of dimension
\(n\), with local coordinates \(q^i\). The cotangent bundle \(T^*Q\) has
induced coordinates \((q^i,p_i)\). When a Lagrangian is regular, the Legendre
map carries \((q,\dot q)\) to \((q,p)\), and the Hamiltonian \(H\) is the
corresponding energy function written in these phase-space coordinates. The
assumptions behind this passage are stated explicitly below, because singular
Lagrangians and gauge constraints require extra reduction machinery that
belongs to more advanced treatments.

\paragraph{Roadmap.}
The chapter has four layers. First, the Legendre transform explains when
velocity data on \(TQ\) can be replaced by momentum data on \(T^*Q\). Second,
the canonical one-form and symplectic form show why Hamilton's equations are
geometric rather than a mnemonic for signs. Third, Poisson brackets, the
phase-space action, and canonical transformations identify the maps that
preserve Hamiltonian structure. Fourth, the numerical lesson explains why a
simulation should preserve geometry as well as approximate trajectories.

Keep one distinction in view throughout the chapter. A point of phase space is
instantaneous data, while a Hamiltonian flow is the rule that moves that data.
The symplectic form is not another force law; it is the structure that turns
the differential of \(H\) into a vector field and decides which coordinate
changes are canonical.

\section{Legendre Transform And Hamiltonian}
\label{sec:legendre-hamiltonian}

The Legendre transform sends \((q,\dot q)\in TQ\) to \((q,p)\in T^*Q\) by
\[
  p_i=\frac{\partial L}{\partial \dot q^i}.
\]
For a regular Lagrangian this can be inverted to express \(\dot q^i\) as a
function of \((q,p,t)\). Concretely, regularity means that the velocity Hessian
\[
  W_{ij}(q,\dot q,t)=
  \frac{\partial^2L}{\partial\dot q^i\partial\dot q^j}
\]
is nonsingular. The inverse function theorem then says that, locally in phase
space, the equations \(p_i=\partial L/\partial\dot q^i\) can be solved for
\(\dot q^i=\dot q^i(q,p,t)\).
The Hamiltonian is
\[
  H(q,p,t)=p_i\dot q^i(q,p,t)-L(q,\dot q(q,p,t),t).
\]

\begin{figure}[htbp]
\centering
\begin{tikzpicture}[scale=1.0, >=Latex, font=\small]
  \node[draw, rounded corners=2pt, minimum width=2.9cm, minimum height=0.9cm,
    fill=blue!5] (tq) at (0,0) {\(TQ:\ (q,\dot q)\)};
  \node[draw, rounded corners=2pt, minimum width=2.9cm, minimum height=0.9cm,
    fill=red!5] (tsq) at (6.2,0) {\(T^*Q:\ (q,p)\)};
  \draw[->, thick] (tq.east) -- node[below=5pt, fill=white, inner sep=2pt]
    {\(\mathbb F L:\ p_i=\partial L/\partial\dot q^i\)} (tsq.west);
  \draw[->, thick]
    ([yshift=5pt]tsq.west) to[out=150,in=30]
    node[above=5pt, fill=white, inner sep=2pt] {regular inverse}
    ([yshift=5pt]tq.east);
  \node[draw, rounded corners=2pt, fill=gray!6, align=center,
    text width=0.46\textwidth, inner sep=4pt] at (3.1,-1.40)
    {regular iff \(\det W\ne 0\)};
\end{tikzpicture}
\caption{The Legendre transform from tangent-bundle data to cotangent-bundle
data. Singular Lagrangians require constraints; the Hamiltonian construction in
this chapter assumes the regular case.}
\label{fig:legendre-transform}
\end{figure}

\begin{derivation}[Hamilton's equations from the Legendre transform]
In this calculation \(H\) is regarded as a function of \((q,p,t)\). The
velocity \(\dot q=\dot q(q,p,t)\) is the inverse image of \((q,p)\) under the
Legendre map, so its differential appears temporarily and then cancels. This is
exactly where regularity is used.
Differentiate \(H\):
\[
\begin{aligned}
  \dd H
  &=
  \dot q^i\,\dd p_i
  +
  p_i\,\dd \dot q^i
  -
  \frac{\partial L}{\partial q^i}\dd q^i
  -
  \frac{\partial L}{\partial \dot q^i}\dd \dot q^i
  -
  \frac{\partial L}{\partial t}\dd t\\
  &=
  \dot q^i\,\dd p_i
  -
  \frac{\partial L}{\partial q^i}\dd q^i
  -
  \frac{\partial L}{\partial t}\dd t,
\end{aligned}
\]
because \(p_i=\partial L/\partial \dot q^i\). Thus
\[
  \frac{\partial H}{\partial p_i}=\dot q^i,
  \qquad
  \frac{\partial H}{\partial q^i}=-\frac{\partial L}{\partial q^i}.
\]
Using the Euler-Lagrange equation \(\dot p_i=\partial L/\partial q^i\), we get
\[
  \boxed{
    \dot q^i=\frac{\partial H}{\partial p_i},
    \qquad
    \dot p_i=-\frac{\partial H}{\partial q^i}
  }.
\]
\end{derivation}

The cancellation of the \(\dd\dot q\) terms is the algebraic heart of the
Legendre transform. It says that once momentum has been defined by
\(p_i=\partial L/\partial\dot q^i\), the Hamiltonian differential contains
only \(\dd q,\dd p,\dd t\). Hamilton's equations are therefore a first-order
system on \(T^*Q\), rather than a second-order system on \(Q\).

The Legendre transform changes the bookkeeping of initial data. A Lagrangian
state \((q,\dot q)\) records a position and a velocity; a Hamiltonian state
\((q,p)\) records a position and the covector that measures how the action
changes under endpoint displacement. The symplectic form below is built from
that endpoint one-form, so Hamiltonian mechanics remembers the variational
origin of the momentum rather than merely renaming velocity.

This is also the point where singular systems leave the main road. If the
velocity Hessian has null directions, some velocities cannot be solved from
the momenta and some momenta become constraints. Electromagnetic gauge theory,
relativistic reparametrization-invariant particles, and ideal incompressible
fluids all have versions of this phenomenon. The present chapter stays in the
regular case so that the geometry of \(T^*Q\) can be learned without the extra
Dirac-constraint machinery.

\section{Phase Space As The Space Of Solutions}
\label{sec:phase-space-solutions}

A useful slogan is that, for a deterministic Hamiltonian system, phase space is
the same thing as the space of solutions.
More precisely, if \(H\) is smooth and the Hamiltonian vector field is complete
on the time interval under study, then every point
\[
  z_0=(q_0,p_0)\in T^*Q
\]
determines a unique trajectory \(z(t)=\Phi^t_H(z_0)\), and every trajectory is
obtained this way from its value at any one time. Thus phase space is the set
of all possible instantaneous states, but it is also a coordinate system on the
family of all classical histories. The caveat is that this identification uses
existence and uniqueness of the equations of motion; singular potentials,
collisions, constraints, or finite-time blowup can invalidate the global
statement even when the local Hamiltonian equations make sense.

\section{The Symplectic Form}
\label{sec:symplectic-form}

The cotangent bundle \(T^*Q\) carries a canonical one-form
\[
  \Theta=p_i\,\dd q^i.
\]
The canonical symplectic form is
\[
  \omega=-\dd\Theta=\dd q^i\wedge\dd p_i.
\]
In arbitrary local phase-space coordinates
\[
  x^m,\qquad m=1,\ldots,2n,
\]
one can write
\[
  \omega=\frac12\omega_{mn}(x)\,\dd x^m\wedge\dd x^n,
\]
where \(\omega_{mn}=-\omega_{nm}\), \(\det(\omega_{mn})\neq0\), and
\[
  \dd\omega=0.
\]
The closedness condition is what makes \(\omega\) locally equivalent to the
canonical form, while nondegeneracy is what lets the Hamiltonian determine a
unique vector field. In canonical coordinates \(x=(q,p)\),
\[
  (\omega_{mn})=
  \begin{pmatrix}
  0 & I_n\\
  -I_n & 0
  \end{pmatrix}.
\]
The Poisson bracket can be expressed intrinsically by the bivector inverse to
the symplectic form, with the sign convention fixed by Hamilton's equations
below. The point is that the bracket is not extra structure once \(\omega\) is
given.

The two defining properties of \(\omega\) have different jobs. Nondegeneracy
lets one solve \(\iota_{X_H}\omega=\dd H\) for a unique vector field, just as a
mass matrix lets one solve force equals mass times acceleration. Closedness is
the integrability condition behind canonical coordinates, Stokes-type
arguments, and the preservation of symplectic area. Hamiltonian mechanics needs
both.

The Hamiltonian vector field \(X_H\) is defined by
\[
  \iota_{X_H}\omega=\dd H.
\]
Here \(\iota_X\) denotes interior contraction: if \(\omega\) is a two-form, then
\((\iota_X\omega)(Y)=\omega(X,Y)\) for every vector field \(Y\).
If
\[
  X_H=A^i\frac{\partial}{\partial q^i}
  +
  B_i\frac{\partial}{\partial p_i},
\]
then
\[
  \iota_{X_H}\omega=A^i\,\dd p_i-B_i\,\dd q^i.
\]
Equating this with
\[
  \dd H=\frac{\partial H}{\partial q^i}\dd q^i
  +
  \frac{\partial H}{\partial p_i}\dd p_i
\]
gives
\[
  A^i=\frac{\partial H}{\partial p_i},
  \qquad
  B_i=-\frac{\partial H}{\partial q^i}.
\]
Thus \(X_H\) is Hamilton's equation written geometrically.

\begin{figure}[htbp]
\centering
\begin{tikzpicture}[scale=1.0, >=Latex]
  \draw[->] (-0.25,0) -- (6.9,0) node[right] {\(q\)};
  \draw[->] (3.35,-2.85) -- (3.35,2.85) node[above] {\(p\)};
  \node[below] at (0.25,0) {\(-\pi\)};
  \node[below] at (3.35,0) {\(0\)};
  \node[below] at (6.45,0) {\(\pi\)};
  \foreach \E/\qmax in {0.35/0.86,0.90/1.47,1.50/2.09}
  {
    \draw[blue!70!black, thick, domain=-\qmax:\qmax, samples=120]
      plot ({\x+3.35},{sqrt(2*(\E-1+cos(\x r)))});
    \draw[blue!70!black, thick, domain=-\qmax:\qmax, samples=120]
      plot ({\x+3.35},{-sqrt(2*(\E-1+cos(\x r)))});
  }
  \draw[red!70!black, thick, domain=-3.12:3.12, samples=180]
    plot ({\x+3.35},{2.0*cos(0.5*\x r)});
  \draw[red!70!black, thick, domain=-3.12:3.12, samples=180]
    plot ({\x+3.35},{-2.0*cos(0.5*\x r)});
  \draw[gray!65, thick, domain=-3.12:3.12, samples=180]
    plot ({\x+3.35},{sqrt(2*(2.35-1+cos(\x r)))});
  \draw[gray!65, thick, domain=-3.12:3.12, samples=180]
    plot ({\x+3.35},{-sqrt(2*(2.35-1+cos(\x r)))});
  \draw[->, thick] (4.15,0.78) -- (4.45,0.62) node[right] {\(X_H\)};
  \node[blue!70!black, fill=white, inner sep=1pt] at (1.30,1.95)
    {libration levels};
  \node[red!70!black, fill=white, inner sep=1pt] at (5.25,-2.10)
    {separatrix};
  \node[gray!65!black, fill=white, inner sep=1pt] at (1.25,-2.35)
    {rotation levels};
  \node at (3.35,-3.18) {Exact level sets of \(H(q,p)=\frac12p^2+1-\cos q\).};
\end{tikzpicture}
\caption{The pendulum phase plane. Libration curves have
\(0<E<2\), the separatrix has \(E=2\), and rotation curves have \(E>2\).
The Hamiltonian vector field is tangent to energy level sets because
\(\dot H=\{H,H\}=0\).}
\label{fig:pendulum-phase-plane}
\end{figure}

\section{Poisson Brackets And Time Evolution}
\label{sec:poisson-time-evolution}

For \(f,g\in C^\infty(T^*Q)\), define
\[
  \Poisson{f}{g}
  =
  \frac{\partial f}{\partial q^i}
  \frac{\partial g}{\partial p_i}
  -
  \frac{\partial f}{\partial p_i}
  \frac{\partial g}{\partial q^i}.
\]
If \(z(t)=(q(t),p(t))\) follows Hamilton's equation, then
\[
  \frac{\dd}{\dd t}f(q(t),p(t),t)
  =
  \frac{\partial f}{\partial t}+\Poisson{f}{H}.
\]
A time-independent observable \(f\) is conserved exactly when
\[
  \Poisson{f}{H}=0.
\]

\begin{proposition}[Liouville volume preservation]
Hamiltonian flow preserves the phase-space volume form
\[
  \Omega=\frac{1}{n!}\omega^n.
\]
\end{proposition}

\begin{proof}
Cartan's formula gives
\[
  \Lie_{X_H}\omega=\dd(\iota_{X_H}\omega)+\iota_{X_H}\dd\omega.
\]
Since \(\iota_{X_H}\omega=\dd H\) and \(\dd\omega=0\),
\[
  \Lie_{X_H}\omega=\dd\dd H=0.
\]
Therefore \(\Lie_{X_H}\omega^n=0\), so the associated volume form is preserved.
\end{proof}

\section{Phase-Space Action}
\label{sec:phase-space-action}

Hamilton's equations can also be derived from the phase-space action
\[
  S_H[q,p]=\int_{t_0}^{t_1}\left(p_i\dot q^i-H(q,p,t)\right)\dd t.
\]
Varying \(q\) and \(p\) independently, with \(\delta q(t_0)=\delta q(t_1)=0\),
gives
\[
  \delta S_H
  =
  \int_{t_0}^{t_1}
  \left[
    \left(\dot q^i-\frac{\partial H}{\partial p_i}\right)\delta p_i
    -
    \left(\dot p_i+\frac{\partial H}{\partial q^i}\right)\delta q^i
  \right]\dd t.
\]
Since \(\delta q\) and \(\delta p\) are arbitrary in the interior, the
stationary curves are exactly Hamiltonian trajectories.

\section{Canonical Transformations}
\label{sec:canonical-transformations}

A diffeomorphism \(\Phi:T^*Q\to T^*Q\) is canonical, or symplectic, if
\[
  \Phi^*\omega=\omega.
\]
Canonical transformations preserve Poisson brackets:
\[
  \Poisson{f\circ\Phi}{g\circ\Phi}
  =
  \Poisson{f}{g}\circ\Phi.
\]
This is why canonical coordinates may change but Hamiltonian mechanics keeps
the same form. Later, action-angle variables will be canonical coordinates
adapted to integrable motion.

\begin{derivation}[Generating function for a canonical transformation]
Let \((q,p)\) be old canonical coordinates and \((Q,P)\) new ones. A convenient
type of generating function depends on the new momenta and old positions:
\[
  F=F(P,q).
\]
Impose
\[
  Q^i=\frac{\partial F}{\partial P_i}\bigg|_q,
  \qquad
  p_i=\frac{\partial F}{\partial q^i}\bigg|_P.
\]
Then
\[
  \dd F=Q^i\,\dd P_i+p_i\,\dd q^i.
\]
Taking another exterior derivative gives
\[
  0=\dd Q^i\wedge\dd P_i+\dd p_i\wedge\dd q^i.
\]
Therefore
\[
  \dd Q^i\wedge\dd P_i
  =
  \dd q^i\wedge\dd p_i.
\]
The canonical symplectic form is preserved, so the transformation is canonical.
In coordinates this is the same cancellation one sees by expanding
\[
  \dd Q^i\wedge\dd P_i
  =
  \left(
    \frac{\partial^2F}{\partial P_i\partial q^j}\dd q^j
    +
    \frac{\partial^2F}{\partial P_i\partial P_j}\dd P_j
  \right)
  \wedge \dd P_i,
\]
where the pure \(\dd P_i\wedge\dd P_j\) term vanishes because second
derivatives are symmetric while the wedge product is antisymmetric.
\end{derivation}

The generating function is useful because it replaces the condition
\(\Phi^*\omega=\omega\), which is a two-form equation, by the exactness of a
one-form. Locally, a canonical transformation is often found by guessing the
right exact differential. Globally, the same calculation must be patched across
coordinate charts, which is why action-angle generating functions can be
multivalued even when the symplectic form is perfectly well defined.

\section{Numerical Lesson}
\label{sec:numerical-lesson}

The equations in this chapter are exact statements about smooth systems. A
numerical integrator creates a discrete map
\[
  z_{n+1}=\Psi_{\Delta t}(z_n).
\]
For long-time Hamiltonian problems, whether \(\Psi_{\Delta t}\) preserves a
symplectic form or a modified Hamiltonian can matter more than the local order
of the method. This is why the pendulum, standard map, and asteroid demos are
not merely programming exercises; they are experiments in preserving structure.

The practical lesson is diagnostic. A good mechanics simulation reports not
only a trajectory but also the quantities that the continuous problem says
should be preserved or nearly preserved. Energy drift, angular-momentum drift,
phase-space area, and event criteria are part of the model audit trail. Without
them, a numerical plot can look convincing while solving the wrong mechanical
problem.

\section{What To Take Forward}
\label{sec:hamiltonian-take-forward}

The rest of the notes repeatedly reuse four lessons from this chapter.
\begin{itemize}
\item The configuration space \(Q\) is part of the model. Choosing bad
coordinates can obscure a simple system, but it cannot change the system.
\item Boundary terms are information, not clutter. They identify canonical
momenta, forces, and conserved quantities.
\item Symmetry produces conserved quantities through Noether's theorem, and
those quantities are often the coordinates of a reduced problem.
\item Hamiltonian mechanics is geometry on \(T^*Q\): the symplectic form
determines Hamilton's equations, volume preservation, Poisson brackets, and the
meaning of a canonical numerical method.
\end{itemize}

\section{Chapter-End Laboratories And Exercises}
\label{sec:hamiltonian-chapter-end-labs}

The laboratories below test when a Lagrangian can be converted into a
Hamiltonian system and how the resulting phase portrait organizes motion.

\subsection{A Nonstandard Lagrangian And Regularity}
\label{sec:lab-nonstandard-lagrangian}

This laboratory separates two ideas that are often conflated. A Lagrangian may
produce a familiar-looking equation of motion and still fail to be globally
regular. Regularity is a statement about whether velocities can be solved from
momenta. Without it, the Legendre transform can fold phase space and the
Hamiltonian description must be handled with care.

Consider the one-dimensional Lagrangian
\[
  L(x,\dot x)
  =
  \frac{m^2\dot x^4}{12}
  +
  m\dot x^2V(x)
  -
  V(x)^2.
\]
The momentum is
\[
  p=\frac{\partial L}{\partial\dot x}
  =
  \frac{m^2}{3}\dot x^3+2mV\dot x.
\]
The Euler-Lagrange equation gives
\[
\begin{aligned}
  0
  &=
  \frac{\dd}{\dd t}
  \left(
    \frac{m^2}{3}\dot x^3+2mV\dot x
  \right)
  -
  \left(
    m\dot x^2V'-2VV'
  \right)\\
  &=
  (m\dot x^2+2V)(m\ddot x+V').
\end{aligned}
\]
There are two branches. The Newtonian branch is
\[
  m\ddot x=-V'(x).
\]
The second branch,
\[
  m\dot x^2+2V(x)=0,
\]
is a singular energy-surface condition. The velocity Hessian
\[
  \frac{\partial^2L}{\partial\dot x^2}
  =
  m(m\dot x^2+2V)
\]
vanishes on this second branch. This example is a useful warning: a Lagrangian
can produce familiar equations on a regular branch while being singular on
another branch. Regularity is not decorative; it is what lets the Legendre
transform define a Hamiltonian system.

\begin{figure}[htbp]
\centering
\begin{tikzpicture}[scale=1.0, >=Latex, font=\small]
  \draw[->] (-3.0,0) -- (3.15,0) node[right] {\(\dot x\)};
  \draw[->] (0,-2.2) -- (0,2.25) node[above] {\(p\)};
  \draw[blue!70!black, thick, domain=-2.55:2.55, samples=180]
    plot (\x,{0.33*\x*\x*\x-0.9*\x});
  \draw[dashed, red!75!black] (-0.95,-1.9) -- (-0.95,1.9);
  \draw[dashed, red!75!black] (0.95,-1.9) -- (0.95,1.9);
  \fill[red!75!black] (-0.95,0.57) circle (1.3pt)
    node[above left] {\(W=0\)};
  \fill[red!75!black] (0.95,-0.57) circle (1.3pt)
    node[below right] {\(W=0\)};
  \draw[decorate, decoration={brace, amplitude=5pt}]
    (-0.95,-1.15) -- (0.95,-1.15)
    node[midway, below=9pt, fill=white, inner sep=1.5pt]
    {multi-valued velocity branch};
  \node[blue!70!black, align=center, fill=white, inner sep=2pt] at (2.05,1.45)
    {folded\\Legendre map};
\end{tikzpicture}
\caption{Nonstandard Lagrangian laboratory. For a representative region with
\(V<0\), the momentum-velocity map can fold where the velocity Hessian
\(W=\partial^2L/\partial\dot x^2\) vanishes. At such points the Legendre
transform is not locally invertible, even though the Euler-Lagrange equation has
a familiar Newtonian branch.}
\label{fig:lab-nonstandard-legendre-fold}
\end{figure}
\FloatBarrier

\subsection{Pendulum Phase Space And The Separatrix}
\label{sec:detailed-pendulum}

The simple pendulum is the canonical one-degree-of-freedom Hamiltonian example.
Let \(\theta\) be the angle from the downward vertical. The Lagrangian is
\[
  L=\frac12m\ell^2\dot\theta^2+mg\ell\cos\theta.
\]
The momentum and Hamiltonian are
\[
  p_\theta=m\ell^2\dot\theta,
  \qquad
  H(\theta,p_\theta)
  =
  \frac{p_\theta^2}{2m\ell^2}-mg\ell\cos\theta.
\]
Hamilton's equations are
\[
  \dot\theta=\frac{p_\theta}{m\ell^2},
  \qquad
  \dot p_\theta=-mg\ell\sin\theta.
\]
The equilibria are
\[
  (\theta,p_\theta)=(0,0)
  \quad\text{and}\quad
  (\pi,0)
\]
modulo \(2\pi\). The downward equilibrium is elliptic; the upward equilibrium
is hyperbolic.

The separatrix is the energy level through the hyperbolic point:
\[
  H=mg\ell.
\]
On this level,
\[
  \frac{p_\theta^2}{2m\ell^2}
  =
  mg\ell(1+\cos\theta)
  =
  2mg\ell\cos^2\frac{\theta}{2}.
\]
Therefore
\[
  p_\theta
  =
  \pm2m\ell\sqrt{g\ell}\cos\frac{\theta}{2}
\]
for \(-\pi\leq\theta\leq\pi\). The separatrix divides libration orbits, which
oscillate around \(\theta=0\), from rotation orbits, which wind around the
circle. This picture is the one-degree-of-freedom ancestor of separatrix
splitting in perturbed Hamiltonian systems.

\section{Associated Demos}
\label{sec:hamiltonian-associated-demos}

The associated phase-space demo is
\path{demos/python/hamiltonian_pendulum.py}.
It plots pendulum trajectories and the separatrix in \((\theta,p_\theta)\).
Read it as a check on this chapter's central claim: phase portraits are not
decorative plots of a solution after the fact, but pictures of Hamiltonian
level sets and the symplectic flow on phase space.
